\section{Markov Decision Process Formulation}
\label{sec:mdp}

% Formulations of this problem conventionally give the policy the full
% per-stage decision: which target, if any, each weapon engages. Section~\ref{sec:auction} shows that half of this decision, target assignment given a set of firing weapons, admits a mechanism with a
% provable approximation guarantee and no learning at all.
We define the MDP so that the policy is only ever asked whether each weapon fires this stage. An MDP, $\mathcal{M}$, is represented as a 5-tuple:
\begin{equation}
\mathcal{M} = (\mathcal{S}, \mathcal{A}, \mathcal{T}, R, \gamma)
\end{equation}
where $\mathcal{S}$ is the state space, $\mathcal{A}$ the action space,
$\mathcal{T}$ the transition kernel, $R$ the reward function, and $\gamma$
the discount factor.

\begin{itemize}
\item \textbf{State ($s_t$)}: for every (weapon, target) pair, the
nine quantities
\begin{equation}
\label{eq:state}
s_t = [u_t;\, v_t;\, d_t;\, r_t;\, \alpha_t;\, f_t;\, ws;\, we;\, p] \in \mathbb{R}^9,
\end{equation}
namely weapon readiness $u_t$, remaining target value $v_t$, remaining reload
time $d_t$, remaining time to horizon $r_t$, remaining ammunition $\alpha_t$,
cumulative hits on the target $f_t$, the window bounds $ws$ and $we$, and the
damage rate $p = P_{m,n}$. Section~\ref{sec:encoding} encodes these as a
heterogeneous graph rather than as $M \times N$ independent vectors.

\item \textbf{Action ($a_t$)}: a \emph{binary} per-weapon participation
decision,
\begin{equation}
\label{eq:binary-action}
a_t = (a_{1,t}, \ldots, a_{M,t}) \in \{0,1\}^M, \qquad a_{m,t} = 1 \iff m \text{ fires at } t,
\end{equation}
rather than a choice of target. This is the entire output the policy is
responsible for: $2^M$ joint actions per stage, against the
$\prod_m (|\mathcal{L}_{m,t}|+1)$ a target-selecting formulation would have to choose from.

\item \textbf{Transition}: given $a_t$, the environment (not the policy)
determines the realized engagement. The set of firing weapons
$F_t = \{m : a_{m,t}=1\}$ is passed to a fixed assignment routine
\begin{equation}
\label{eq:auction-call}
x_t = \textsc{Assign}(F_t, s_t),
\end{equation}
which gives each $m \in F_t$ one target, choosing greedily by the value a
shot would destroy at the current state. It is deterministic, has no learned
parameters, and is defined in full in Section~\ref{sec:auction}. A weapon in
$F_t$ with no legal target is a no-op regardless of $a_{m,t}$. The transition
$\mathcal{T}(s_{t+1} \mid s_t, a_t)$ is therefore deterministic once $a_t$
is fixed, but it is not simply ``apply $a_t$'': it composes the policy's
binary decision with the fixed assignment rule.

\item \textbf{Reward}:
\begin{equation}
R(s_{t+1} \mid s_t, a_t) = \frac{C(s_t) - C(s_{t+1})}{C(s_0)},
\end{equation}
where $C(s_t)$ is the sum of remaining target values at $s_t$, realized
through the assigned engagement rather than a directly-specified one.
\end{itemize}

This factoring of the decision, learn participation and delegate assignment,
is what Section~\ref{sec:method} implements, and what
Section~\ref{sec:participation-theory}'s myopic-gap proposition argues cannot be
undone: the assignment rule alone, applied to \emph{every} weapon with a
positive marginal value each stage, is a valid (if unlearned) participation
rule, and
Proposition~\ref{prop:myopic-gap} shows it can be arbitrarily far from
optimal. The gap between that myopic rule and the optimal one is exactly
what Eq.~\eqref{eq:binary-action}'s policy is trained to close.
